% Options for packages loaded elsewhere
\PassOptionsToPackage{unicode}{hyperref}
\PassOptionsToPackage{hyphens}{url}
\documentclass[
]{article}
\usepackage{xcolor}
\usepackage[margin=1in]{geometry}
\usepackage{amsmath,amssymb}
\setcounter{secnumdepth}{5}
\usepackage{iftex}
\ifPDFTeX
  \usepackage[T1]{fontenc}
  \usepackage[utf8]{inputenc}
  \usepackage{textcomp} % provide euro and other symbols
\else % if luatex or xetex
  \usepackage{unicode-math} % this also loads fontspec
  \defaultfontfeatures{Scale=MatchLowercase}
  \defaultfontfeatures[\rmfamily]{Ligatures=TeX,Scale=1}
\fi
\usepackage{lmodern}
\ifPDFTeX\else
  % xetex/luatex font selection
\fi
% Use upquote if available, for straight quotes in verbatim environments
\IfFileExists{upquote.sty}{\usepackage{upquote}}{}
\IfFileExists{microtype.sty}{% use microtype if available
  \usepackage[]{microtype}
  \UseMicrotypeSet[protrusion]{basicmath} % disable protrusion for tt fonts
}{}
\makeatletter
\@ifundefined{KOMAClassName}{% if non-KOMA class
  \IfFileExists{parskip.sty}{%
    \usepackage{parskip}
  }{% else
    \setlength{\parindent}{0pt}
    \setlength{\parskip}{6pt plus 2pt minus 1pt}}
}{% if KOMA class
  \KOMAoptions{parskip=half}}
\makeatother
\usepackage{color}
\usepackage{fancyvrb}
\newcommand{\VerbBar}{|}
\newcommand{\VERB}{\Verb[commandchars=\\\{\}]}
\DefineVerbatimEnvironment{Highlighting}{Verbatim}{commandchars=\\\{\}}
% Add ',fontsize=\small' for more characters per line
\usepackage{framed}
\definecolor{shadecolor}{RGB}{248,248,248}
\newenvironment{Shaded}{\begin{snugshade}}{\end{snugshade}}
\newcommand{\AlertTok}[1]{\textcolor[rgb]{0.94,0.16,0.16}{#1}}
\newcommand{\AnnotationTok}[1]{\textcolor[rgb]{0.56,0.35,0.01}{\textbf{\textit{#1}}}}
\newcommand{\AttributeTok}[1]{\textcolor[rgb]{0.13,0.29,0.53}{#1}}
\newcommand{\BaseNTok}[1]{\textcolor[rgb]{0.00,0.00,0.81}{#1}}
\newcommand{\BuiltInTok}[1]{#1}
\newcommand{\CharTok}[1]{\textcolor[rgb]{0.31,0.60,0.02}{#1}}
\newcommand{\CommentTok}[1]{\textcolor[rgb]{0.56,0.35,0.01}{\textit{#1}}}
\newcommand{\CommentVarTok}[1]{\textcolor[rgb]{0.56,0.35,0.01}{\textbf{\textit{#1}}}}
\newcommand{\ConstantTok}[1]{\textcolor[rgb]{0.56,0.35,0.01}{#1}}
\newcommand{\ControlFlowTok}[1]{\textcolor[rgb]{0.13,0.29,0.53}{\textbf{#1}}}
\newcommand{\DataTypeTok}[1]{\textcolor[rgb]{0.13,0.29,0.53}{#1}}
\newcommand{\DecValTok}[1]{\textcolor[rgb]{0.00,0.00,0.81}{#1}}
\newcommand{\DocumentationTok}[1]{\textcolor[rgb]{0.56,0.35,0.01}{\textbf{\textit{#1}}}}
\newcommand{\ErrorTok}[1]{\textcolor[rgb]{0.64,0.00,0.00}{\textbf{#1}}}
\newcommand{\ExtensionTok}[1]{#1}
\newcommand{\FloatTok}[1]{\textcolor[rgb]{0.00,0.00,0.81}{#1}}
\newcommand{\FunctionTok}[1]{\textcolor[rgb]{0.13,0.29,0.53}{\textbf{#1}}}
\newcommand{\ImportTok}[1]{#1}
\newcommand{\InformationTok}[1]{\textcolor[rgb]{0.56,0.35,0.01}{\textbf{\textit{#1}}}}
\newcommand{\KeywordTok}[1]{\textcolor[rgb]{0.13,0.29,0.53}{\textbf{#1}}}
\newcommand{\NormalTok}[1]{#1}
\newcommand{\OperatorTok}[1]{\textcolor[rgb]{0.81,0.36,0.00}{\textbf{#1}}}
\newcommand{\OtherTok}[1]{\textcolor[rgb]{0.56,0.35,0.01}{#1}}
\newcommand{\PreprocessorTok}[1]{\textcolor[rgb]{0.56,0.35,0.01}{\textit{#1}}}
\newcommand{\RegionMarkerTok}[1]{#1}
\newcommand{\SpecialCharTok}[1]{\textcolor[rgb]{0.81,0.36,0.00}{\textbf{#1}}}
\newcommand{\SpecialStringTok}[1]{\textcolor[rgb]{0.31,0.60,0.02}{#1}}
\newcommand{\StringTok}[1]{\textcolor[rgb]{0.31,0.60,0.02}{#1}}
\newcommand{\VariableTok}[1]{\textcolor[rgb]{0.00,0.00,0.00}{#1}}
\newcommand{\VerbatimStringTok}[1]{\textcolor[rgb]{0.31,0.60,0.02}{#1}}
\newcommand{\WarningTok}[1]{\textcolor[rgb]{0.56,0.35,0.01}{\textbf{\textit{#1}}}}
\usepackage{longtable,booktabs,array}
\newcounter{none} % for unnumbered tables
\usepackage{calc} % for calculating minipage widths
% Correct order of tables after \paragraph or \subparagraph
\usepackage{etoolbox}
\makeatletter
\patchcmd\longtable{\par}{\if@noskipsec\mbox{}\fi\par}{}{}
\makeatother
% Allow footnotes in longtable head/foot
\IfFileExists{footnotehyper.sty}{\usepackage{footnotehyper}}{\usepackage{footnote}}
\makesavenoteenv{longtable}
\usepackage{graphicx}
\makeatletter
\newsavebox\pandoc@box
\newcommand*\pandocbounded[1]{% scales image to fit in text height/width
  \sbox\pandoc@box{#1}%
  \Gscale@div\@tempa{\textheight}{\dimexpr\ht\pandoc@box+\dp\pandoc@box\relax}%
  \Gscale@div\@tempb{\linewidth}{\wd\pandoc@box}%
  \ifdim\@tempb\p@<\@tempa\p@\let\@tempa\@tempb\fi% select the smaller of both
  \ifdim\@tempa\p@<\p@\scalebox{\@tempa}{\usebox\pandoc@box}%
  \else\usebox{\pandoc@box}%
  \fi%
}
% Set default figure placement to htbp
\def\fps@figure{htbp}
\makeatother
% definitions for citeproc citations
\NewDocumentCommand\citeproctext{}{}
\NewDocumentCommand\citeproc{mm}{%
  \begingroup\def\citeproctext{#2}\cite{#1}\endgroup}
\makeatletter
 % allow citations to break across lines
 \let\@cite@ofmt\@firstofone
 % avoid brackets around text for \cite:
 \def\@biblabel#1{}
 \def\@cite#1#2{{#1\if@tempswa , #2\fi}}
\makeatother
\newlength{\cslhangindent}
\setlength{\cslhangindent}{1.5em}
\newlength{\csllabelwidth}
\setlength{\csllabelwidth}{3em}
\newenvironment{CSLReferences}[2] % #1 hanging-indent, #2 entry-spacing
 {\begin{list}{}{%
  \setlength{\itemindent}{0pt}
  \setlength{\leftmargin}{0pt}
  \setlength{\parsep}{0pt}
  % turn on hanging indent if param 1 is 1
  \ifodd #1
   \setlength{\leftmargin}{\cslhangindent}
   \setlength{\itemindent}{-1\cslhangindent}
  \fi
  % set entry spacing
  \setlength{\itemsep}{#2\baselineskip}}}
 {\end{list}}
\usepackage{calc}
\newcommand{\CSLBlock}[1]{\hfill\break\parbox[t]{\linewidth}{\strut\ignorespaces#1\strut}}
\newcommand{\CSLLeftMargin}[1]{\parbox[t]{\csllabelwidth}{\strut#1\strut}}
\newcommand{\CSLRightInline}[1]{\parbox[t]{\linewidth - \csllabelwidth}{\strut#1\strut}}
\newcommand{\CSLIndent}[1]{\hspace{\cslhangindent}#1}
\setlength{\emergencystretch}{3em} % prevent overfull lines
\providecommand{\tightlist}{%
  \setlength{\itemsep}{0pt}\setlength{\parskip}{0pt}}
\usepackage{bookmark}
\IfFileExists{xurl.sty}{\usepackage{xurl}}{} % add URL line breaks if available
\urlstyle{same}
\hypersetup{
  pdftitle={Because we can: A unified Bayesian engine for causal inference across biological scales},
  pdfkeywords={Causal Inference, Bayesian Structural Equation Modeling, Directed Acyclic Graph, Phylogenetic Comparative Methods, Spatial Autocorrelation, Markov chain Monte Carlo},
  hidelinks,
  pdfcreator={LaTeX via pandoc}}

\title{Because we can: A unified Bayesian engine for causal inference across biological scales}
\author{true \and true}
\date{}

\begin{document}
\maketitle
\begin{abstract}
Moving beyond description to explanation is the hallmark of a mature science. As ecology and evolutionary biology increasingly adopt formal causal inference frameworks, researchers are tasked with parameterizing complex causal topologies with empirical data. However, biological data is uniquely characterized by multi-level dependencies, from deep evolutionary history to contemporary spatial autocorrelation. Historically, researchers have been forced into a methodological compromise: use accessible piecewise approaches that struggle to propagate uncertainty across latent states, or build custom, fully joint Bayesian Structural Equation Models (BSEMs) that are technically daunting and computationally crippled by the ``Curse of Covariance.'' Here, we introduce the \texttt{because} (Bayesian Estimation of Causal Effects) R package, an integrated framework designed to serve as a unified engine for causal inference across biological scales. By providing a ``Mechanism-First'' declarative syntax, the framework automates the translation of theoretical Directed Acyclic Graphs (DAGs) into optimized Bayesian statistical implementations. Crucially, \texttt{because} automates the structural re-parameterization required to bypass dense matrix inversions, effectively linearizing the computational cost of spatial and phylogenetic models. We detail the architecture of this engine, its handling of non-Gaussian likelihoods and latent states, and present case studies demonstrating its ability to empower researchers to focus on the structure of their systems rather than the computational limits of their software.
\end{abstract}

\section{1. Introduction: Bridging Causal Theory Across Biological Scales}\label{introduction-bridging-causal-theory-across-biological-scales}

As the ecological and evolutionary sciences increasingly adopt formal causal inference frameworks, researchers are no longer solely asking \emph{whether} to employ Directed Acyclic Graphs (DAGs), but \emph{how} to robustly parameterize them with empirical data (Siegel and Dee 2025; Dudney et al. 2025). A defining feature of biology is that causality operates simultaneously across multiple scales of space and time. A species' functional traits are shaped by deep evolutionary history (Felsenstein 1985), while its local abundance is constrained by contemporary spatial autocorrelation (Legendre 1993) and hierarchical population structures. A complete mechanistic explanation of biological phenomena requires a framework capable of disentangling these overlapping dependencies within a unified causal topology.

Historically, the tools available for causal inference have enforced an artificial segregation of these scales. Evolutionary biologists often rely on specialized Phylogenetic Comparative Methods (PCMs; e.g., Hadfield (2010)) to detect trait correlations across large, non-independent clades. Conversely, ecological research has widely adopted path analysis or piecewise Structural Equation Modeling (piecewiseSEM; Lefcheck (2016)) to estimate local, hierarchical causal effects using familiar linear mixed-model frameworks. The challenge arises when research demands the integration of both. For example, quantifying how climate change (an ecological exposure) indirectly alters species distributions via its effect on thermal tolerance (an evolutionarily constrained mediator) requires a causal model that simultaneously accounts for both local spatial diffusion and deep phylogenetic inertia within the same network.

Presently, uniting these scales requires researchers to abandon ``off-the-shelf'' software and construct custom, fully joint Bayesian Structural Equation Models (BSEMs) in languages like JAGS or Stan. While analytically powerful (Hardenberg and Gonzalez-Voyer 2025), custom-scripting these models demands specialized statistical programming expertise. Furthermore, such joint models suffer from a severe computational bottleneck. In standard model formulations, Markov chain Monte Carlo (MCMC) algorithms must repeatedly invert massive, multi-scale covariance matrices at every iteration---an operation that scales cubically (\(O(N^3)\)) with dataset size, rendering large networks computationally prohibitive.

While the computational advantages of conditional hierarchical parameterization and pre-computed precision matrices have been successfully leveraged to solve this bottleneck in specialized single-response mixed-model software (e.g., MCMCglmm, Hadfield (2010); INLA, Rue et al. (2009)), extending these algebraic optimizations to complex, multi-equation causal networks has historically required prohibitive custom programming.

This paper introduces \texttt{because} (Bayesian Estimation of Causal Effects), an integrated R package designed to serve as a unified engine for causal inference across biological scales. The central methodological contribution of \texttt{because} is its ability to seamlessly bridge evolutionary and ecological hierarchies within a single generative Bayesian framework. It provides researchers with a ``Mechanism-First'' declarative syntax, automating the complex translation of theoretical DAGs into fully joint hierarchical models. Crucially, \texttt{because} automates the structural re-parameterization required to bypass dense matrix inversions across the entire causal graph, democratizing highly optimized, joint Bayesian inference for the life sciences.

In this contribution, we detail the architecture and capabilities of this engine. In \textbf{Section 2}, we outline the declarative framework that allows researchers to specify complex, multi-scale topologies, demonstrating how standard R formulas are assembled into robust causal models. In \textbf{Section 3}, we detail the computational infrastructure that handles generalized covariance structures---from spatial autocorrelation to multi-tree phylogenetic uncertainty---allowing evolutionary and ecological processes to be evaluated simultaneously. In \textbf{Section 4}, we explore how the framework's native handling of latent states allows for the rigorous causal treatment of missing data and non-detection. Finally, in \textbf{Section 5}, we present case studies demonstrating the framework's ability to disentangle complex, multi-scale biological phenomena, highlighting how \texttt{because} empowers researchers to focus on the structure of their systems rather than the computational limits of their software.

\section{2. Framework \& Implementation: The Mechanism-First Workflow}\label{framework-implementation-the-mechanism-first-workflow}

The primary barrier to widespread adoption of joint Bayesian structural equation models in ecology and evolution is the steep technical curve required to specify them. Writing custom Markov chain Monte Carlo (MCMC) samplers requires researchers to explicitly define every prior, manage complex hierarchical distributions, and manually construct likelihood functions for non-standard data types. This technical burden often forces researchers into a ``Software-First'' workflow, where the structure of the biological hypothesis is constrained by the limitations of the most accessible software package.

\texttt{because} inverts this paradigm, enabling a ``Mechanism-First'' workflow. The core philosophy of the framework is that the primary task of the biologist is the formal declaration of the theoretical Directed Acyclic Graph (DAG; Pearl (2009)). Once the biological topology is declared, the software assumes the burden of statistical implementation.

\subsection{2.1. Declarative Model Specification}\label{declarative-model-specification}

In \texttt{because}, the causal network is specified using a list of standard \texttt{R} formulas. Each formula represents a distinct node in the DAG and its direct parents (causes). For example, a system where environmental temperature (\texttt{Temp}) drives both body mass (\texttt{Mass}) and, indirectly, overwinter survival (\texttt{Surv}), is declared simply as:

\begin{Shaded}
\begin{Highlighting}[]
\NormalTok{eqs }\OtherTok{\textless{}{-}} \FunctionTok{list}\NormalTok{(}
\NormalTok{  Mass }\SpecialCharTok{\textasciitilde{}}\NormalTok{ Temp,}
\NormalTok{  Surv }\SpecialCharTok{\textasciitilde{}}\NormalTok{ Mass }\SpecialCharTok{+}\NormalTok{ Temp}
\NormalTok{)}
\end{Highlighting}
\end{Shaded}

This declarative syntax is deliberately minimalistic. Behind the scenes, \texttt{because} parses this list to construct the full causal topology. By analyzing the structural equations, the engine automatically resolves the implied d-separation statements (the set of conditional independencies that must hold true if the DAG represents the true data-generating process). These tests are generated dynamically, allowing the researcher to immediately validate their assumed causal structure against the empirical data without writing custom testing scripts (Shipley 2000, 2013).

\subsection{2.2. Automating the Structural Translation}\label{automating-the-structural-translation}

The mathematical translation of the DAG into a Bayesian framework is where the \texttt{because} engine provides its most significant automation. Unlike piecewise frequentist approaches that estimate each equation independently, \texttt{because} compiles the entire system of equations into a single, joint JAGS model string (Plummer 2003).

This automated compilation handles several critical components of causal modeling that are notoriously error-prone when coded manually:

\begin{enumerate}
\def\labelenumi{\arabic{enumi}.}
\tightlist
\item
  \textbf{Prior Generation:} The engine automatically assigns appropriate, weakly informative priors to all intercept and coefficient terms, scaling them relative to the variance of the data to ensure standardized MCMC behavior across disparate datasets.
\item
  \textbf{Multivariate Design Matrices:} Categorical confounders and interactive effects are automatically dummy-coded and mapped to their appropriate hierarchical levels within the JAGS syntax.
\item
  \textbf{Missing Data Imputation:} In a joint Bayesian framework, missing data is not an error; it is a latent parameter. \texttt{because} automatically identifies missing observations (\texttt{NA} values) across the dataset and constructs the necessary stochastic nodes to perform MCMC-based imputation, ensuring that uncertainty from unobserved states is propagated through the entire causal network.
\end{enumerate}

\subsection{2.3. Deterministic Nodes and Biological Logic}\label{deterministic-nodes-and-biological-logic}

Biological mechanisms are rarely purely linear or additive. They frequently involve hard thresholds, nonlinear transformations, or logical states. To accommodate this, \texttt{because} supports the inclusion of \emph{deterministic nodes} within the causal graph.

Researchers can embed logical statements directly into the structural equations (e.g., \texttt{Surv\ \textasciitilde{}\ Mass\ +\ I(Temp\ \textgreater{}\ 30)} or \texttt{Fecundity\ \textasciitilde{}\ I(Mass\^{}2)}). The engine parses these standard \texttt{R} expressions and translates them into deterministic step functions (\texttt{step()}) or algebraic transformations within the JAGS model. This capability is vital for ecological causal inference, as it allows the model to faithfully reflect non-linear biological realities---such as critical thermal maximums or threshold carrying capacities---without forcing the user to step outside the bounds of the declarative network syntax.

\section{\texorpdfstring{3. The \texttt{because} Engine: Overcoming the Covariance Bottleneck}{3. The because Engine: Overcoming the Covariance Bottleneck}}\label{the-because-engine-overcoming-the-covariance-bottleneck}

The ability to easily declare a complex causal topology is only half of the ``Mechanism-First'' paradigm. The other half is the ability to computationally fit that topology to empirical data within a reasonable timeframe. It is here that custom-coded Bayesian Structural Equation Models have historically stumbled, falling victim to the computational realities of non-independent data.

\subsection{3.1. The Curse of Covariance}\label{the-curse-of-covariance}

When biological data exhibits overlapping dependencies---such as phenotypic traits constrained by a shared phylogenetic tree and sampled across a spatially autocorrelated landscape---these relationships must be explicitly modeled to avoid spurious causal inference. In traditional Bayesian formulations, these dependencies are often handled marginally via a multivariate normal likelihood (\(Y \sim \mathcal{N}(\mu, \Sigma)\)), where the covariance matrix \(\Sigma\) is a function of the known structure (e.g., the phylogenetic variance-covariance matrix, \(\mathbf{V}\)) scaled by an estimated variance parameter (\(\sigma^2\)).

Because the scaling parameter (\(\sigma^2\)) is updated at every step of the MCMC algorithm, the sampler must recalculate and invert the entire dense \(N \times N\) covariance matrix at every single iteration. Matrix inversion is an operation with a computational complexity of \(O(N^3)\). Consequently, what is computationally trivial for 50 taxa becomes paralyzing for 500 taxa, placing a hard limit on the scale and complexity of BSEMs that can be practically evaluated.

\subsection{3.2. Structural Re-parameterization}\label{structural-re-parameterization}

\texttt{because} circumvents this bottleneck by automating a structural re-parameterization of the causal graph, shifting from a marginal likelihood to a conditional hierarchical formulation. While this mathematical trick has been successfully deployed in quantitative genetics (Henderson 1950) and specialized single-response mixed-model engines (e.g., \texttt{MCMCglmm}; Hadfield (2010)), extending it across a multi-node, directed causal network has historically required prohibitive custom matrix algebra.

Instead of modeling the structured variance marginally, \texttt{because} separates the structural (causal) expectations from the stochastic (covariance) expectations. The response variable is modeled as conditionally independent (\(Y_i \sim \mathcal{N}(\mu_i + u_i, \sigma^2_e)\)), where \(u_i\) is a latent parameter capturing the structured residual. The structured residuals themselves are then modeled as a multivariate normal distribution centered on zero (\(u \sim \mathcal{N}(0, \lambda \mathbf{P})\)).

The critical computational advantage lies in the term \(\mathbf{P}\). Rather than repeatedly inverting a covariance matrix (\(\mathbf{V}\)) inside the MCMC sampler, \texttt{because} analytically calculates the \emph{precision matrix} (\(\mathbf{P} = \mathbf{V}^{-1}\)) a single time in \texttt{R} prior to model compilation. This fixed precision matrix is then passed to JAGS as data. During MCMC sampling, the computationally expensive \(O(N^3)\) matrix inversion is replaced by a simple \(O(N^2)\) scalar-matrix multiplication (\(\lambda \cdot \mathbf{P}\)).

\subsection{3.3. Unifying Evolutionary and Ecological Scales}\label{unifying-evolutionary-and-ecological-scales}

This automated structural re-parameterization does more than just accelerate model fitting (often resulting in 4-10x speedups over marginal implementations); it serves as the enabling technology that allows \texttt{because} to unify disparate biological scales.

Because the \(O(N^3)\) computational penalty has been removed, researchers are no longer forced to choose between modeling deep evolutionary history or contemporary local ecology. \texttt{because} allows for the simultaneous inclusion of precision matrices derived from phylogenetic trees, spatial distance-decay networks, and hierarchical grouping factors (e.g., generic random intercepts) within the same causal graph. Furthermore, this efficiency allows the engine to seamlessly handle phylogenetic uncertainty---sampling across a posterior distribution of hundreds of precision matrices during a single MCMC run---without stalling the sampler.

By automating the algebra required to bypass the covariance bottleneck, \texttt{because} provides researchers with an engine capable of testing causal hypotheses at the true, multi-scale complexity of the biological systems they study.

\section{4. Core Features: Expanding the Causal Toolkit}\label{core-features-expanding-the-causal-toolkit}

The computational architecture outlined in Section 3 provides the foundation for \texttt{because} to tackle the realities of ecological and evolutionary datasets. In this section, we detail the core extensions supported by the framework, providing the mathematical implementation strategies that allow the engine to maintain Bayesian rigor across diverse data structures.

\subsection{4.1. Non-Gaussian Likelihoods and Link Functions}\label{non-gaussian-likelihoods-and-link-functions}

Biological data is rarely normally distributed. Counts of species, binary survival outcomes, and ordinal behavioral states are ubiquitous. While frequentist Generalized Linear Mixed Models (GLMMs) handle these distributions routinely, incorporating them into a joint causal network requires careful parameterization.

\texttt{because} supports a wide array of non-Gaussian families (e.g., binomial, Poisson, negative binomial, zero-inflated distributions (Lambert 1992), ordinal logistic). When a non-Gaussian family is specified for a node \(Y\), the framework automatically implements the appropriate link function \(g(\cdot)\) and introduces an observation-level structured residual \(e_i\) to absorb overdispersion that would otherwise distort the causal estimates:

\[ g(\mu_i) = \alpha + \sum_{j=1}^{K} \beta_j X_{ji} + e_i \]
\[ Y_i \sim \mathcal{F}(\mu_i, \theta) \]

where \(X_{ji}\) are the \(K\) parent nodes (causes) of \(Y\), \(e_i \sim \mathcal{N}(0, \sigma^2_e)\) is the observation-level residual, and \(\mathcal{F}\) is the chosen distributional family with potential dispersion parameter \(\theta\). For example, in a zero-inflated Poisson (ZIP) model, the network automatically forks the node evaluation into a structural equation for the count process (\(\mu_i\)) and a parallel logit equation for the structural zero-inflation probability (\(\psi_i\)), allowing the user to map causal paths independently to occurrence and abundance.

\subsection{4.2. Missing Data as a Latent Process}\label{missing-data-as-a-latent-process}

A major philosophical and practical limitation of many statistical frameworks is the requirement for complete case data. Listwise deletion (dropping any row with an \texttt{NA} value) not only reduces statistical power but, within a causal network, can introduce severe systemic biases and break causal mediation chains. In a Bayesian paradigm, missing data is not an error; it is an unobserved latent state (Little and Rubin 2019).

\texttt{because} treats all structural equations as generative processes. If a variable \(M\) is partially observed and acts as a mediator between an exposure \(X\) and an outcome \(Y\) (\(X \rightarrow M \rightarrow Y\)), the framework evaluates the known values of \(M\) based on its parents:

\[ M_i^{obs} \sim \mathcal{N}(\alpha_M + \beta_{XM} X_i, \sigma^2_M) \]

For any iteration where \(M_i\) is missing (\(M_i^{miss}\)), JAGS treats it as a stochastic node to be estimated. Because the model is evaluated jointly, the posterior distribution of \(M_i^{miss}\) is constrained simultaneously by its theoretical causes (\(X\)) and its downstream effects on \(Y\). This allows \texttt{because} to perform full-information imputation during model fitting, preserving the causal integrity of the network and accurately propagating the uncertainty of the missing state into the final coefficient estimates.

\subsection{\texorpdfstring{4.3. Extensibility and the \texttt{because.phybase} Module}{4.3. Extensibility and the because.phybase Module}}\label{extensibility-and-the-because.phybase-module}

To ensure that the framework remains modular and adaptable to future methodological advances, \texttt{because} employs an S3-based extension system. This allows researchers to define custom covariance structures and link them into the causal graph without modifying the core engine.

The most prominent implementation of this extensibility is the \texttt{because.phybase} module, which specializes the engine for Phylogenetic Comparative Methods. As discussed in Section 3, incorporating evolutionary history requires evaluating the precision matrix of the phylogenetic tree. When a user provides a \texttt{phylo} object to the model, \texttt{because.phybase} intercepts the call, standardizes the branch lengths, and calculates the precision matrix \(\mathbf{P}\). It then dynamically injects the following nested structure into the base JAGS code:

\[ u_i \sim \mathcal{N}(0, \lambda \mathbf{P}) \]
\[ \lambda \sim \text{Gamma}(0.001, 0.001) \]

Here, \(\lambda\) acts as the phylogenetic scaling parameter (analogous to Pagel's \(\lambda\) in a residual framework). If the researcher passes a \texttt{multiPhylo} object to represent a posterior distribution of trees (incorporating phylogenetic uncertainty), the module automatically upgrades the precision matrix to a 3-dimensional array (\(\mathbf{P}_{1:N, 1:N, 1:K}\)). During MCMC sampling, the model treats the tree index \(k\) as a categorical random variable with a uniform prior, stepping randomly across the \(K\) available precision matrices at each iteration. This allows the model to simultaneously integrate over both parameter space and topological tree space.

\subsection{4.4 Cross-Level Hierarchical Causal Inference}\label{cross-level-hierarchical-causal-inference}

A pervasive challenge in ecology and evolution is that variables frequently operate at different scales of biological organization. For example, a researcher may wish to model how site-level meteorological conditions constrain species-level functional traits, which in turn drive observation-level local abundances. When using piece-wise methods, evaluating causal effects across these entirely different hierarchical levels often forces researchers to improperly aggregate their data (losing valuable variance and statistical power) or manually specify incredibly complex mixed-effect correlation structures.

Because the \texttt{because} engine compiles the full joint distribution into a hierarchical JAGS template, these cross-level causal relationships are evaluated natively. Predictors and responses can exist at entirely different levels of aggregation, and the engine automatically handles the necessary hierarchical indexing under the hood to preserve the precise variance structure at each nested scale.

\section{5. Case Study: Disentangling Evolutionary and Spatial Scales}\label{case-study-disentangling-evolutionary-and-spatial-scales}

To demonstrate the ``Mechanism-First'' workflow and the power of the \texttt{because} engine, we present a simulated case study that explicitly tackles the challenge of dual-scale biological covariance.

\subsubsection{5.1 Simulated Data Generation}\label{simulated-data-generation}

To provide a rigorous test of the joint estimation architecture, we simulated a cross-classified macroecological dataset representing \(N=500\) standardized surveys (e.g., pitfall traps) conducted across \(M=50\) geographically distinct sites. During these surveys, researchers recorded the abundance of \(S=50\) possible species. Crucially, this dataset contains variables operating at three distinct hierarchical levels of biological organization: \texttt{Elevation} and \texttt{Temperature} are site-level properties, \texttt{Body\_Mass} and \texttt{Thermal\_Tolerance} are species-specific evolutionary traits, and \texttt{Abundance} is an observation-level outcome.

We generated the true biological relationships according to a cross-level mediation hypothesis: \texttt{Elevation} drives local \texttt{Temperature} (site-level effect) while a species' \texttt{Body\_Mass} constrains its \texttt{Thermal\_Tolerance} (species-level effect). Finally, \texttt{Temperature} and \texttt{Thermal\_Tolerance} converge to determine local \texttt{Abundance} (observation-level effect), where larger mismatches between the site temperature and species tolerance severely limit population sizes.

Crucially, we injected non-independence at two distinct scales into the data generation process (full \texttt{R} code is provided in the Supplementary Materials):
1. \textbf{Evolutionary Scale:} \texttt{Thermal\_Tolerance} was simulated with strong phylogenetic signal (\(\lambda_{true} = 0.7\)) by generating a 50-tip coalescent tree and drawing baseline tolerance values from a multivariate normal distribution defined by the tree's variance-covariance matrix.
2. \textbf{Ecological Scale:} We assigned random \((X,Y)\) coordinates to the 50 sites and generated a spatial correlation matrix using an exponential distance-decay function (\(W = e^{-0.05 \times D_{ij}}\)). Site-level variance components were then drawn from a multivariate normal distribution based on this spatial matrix, Mathematically encoding Tobler's First Law of Geography into the local temperature baselines.

\subsubsection{5.2 Model Specification and Recovery}\label{model-specification-and-recovery}

While the phylogenetic structure is easily handled by passing the tree via the \texttt{because.phybase} extension, we can easily construct a custom spatial covariance matrix using the \texttt{because\_structure()} helper function. This function allows researchers to extend the package without needing to write complex S3 methods or JAGS code. We simply define the mathematics for generating the precision matrix:

\begin{Shaded}
\begin{Highlighting}[]
\CommentTok{\# Define a custom spatial decay structure}
\NormalTok{spatial\_decay }\OtherTok{\textless{}{-}} \FunctionTok{because\_structure}\NormalTok{(}
    \AttributeTok{name =} \StringTok{"spatial\_decay"}\NormalTok{,}
    \AttributeTok{description =} \StringTok{"Exponential distance decay spatial structure"}\NormalTok{,}
    \AttributeTok{precision\_fn =} \ControlFlowTok{function}\NormalTok{(coords, }\AttributeTok{decay\_rate =} \FloatTok{0.05}\NormalTok{) \{}
        \CommentTok{\# Calculate distance matrix from coordinates}
\NormalTok{        dist\_mat }\OtherTok{\textless{}{-}} \FunctionTok{as.matrix}\NormalTok{(}\FunctionTok{dist}\NormalTok{(coords))}
        \CommentTok{\# Create exponential decay covariance matrix}
\NormalTok{        W }\OtherTok{\textless{}{-}} \FunctionTok{exp}\NormalTok{(}\SpecialCharTok{{-}}\NormalTok{decay\_rate }\SpecialCharTok{*}\NormalTok{ dist\_mat)}
        \CommentTok{\# Add a small ridge for numerical stability}
        \FunctionTok{diag}\NormalTok{(W) }\OtherTok{\textless{}{-}} \FunctionTok{diag}\NormalTok{(W) }\SpecialCharTok{+} \FloatTok{1e{-}4}
        \CommentTok{\# Return the precision matrix (V\^{}{-}1)}
        \FunctionTok{solve}\NormalTok{(W) }
\NormalTok{    \}}
\NormalTok{)}

\CommentTok{\# Instantiate the structure with our site coordinates}
\NormalTok{my\_spatial }\OtherTok{\textless{}{-}} \FunctionTok{spatial\_decay}\NormalTok{(site\_coords)}
\end{Highlighting}
\end{Shaded}

With the evolutionary and ecological structures defined, the researcher can declare the biological hypothesis using standard \texttt{R} formulas and fit the joint Bayesian model:

\begin{Shaded}
\begin{Highlighting}[]
\CommentTok{\# 1. Declare the multi{-}level causal network}
\NormalTok{eqs }\OtherTok{\textless{}{-}} \FunctionTok{list}\NormalTok{(}
\NormalTok{    Thermal\_Tolerance }\SpecialCharTok{\textasciitilde{}}\NormalTok{ Body\_Mass\_s,}
\NormalTok{    Temperature }\SpecialCharTok{\textasciitilde{}}\NormalTok{ Elevation\_s }\SpecialCharTok{+}\NormalTok{ (}\DecValTok{1} \SpecialCharTok{|}\NormalTok{ Site),}
\NormalTok{    Abundance }\SpecialCharTok{\textasciitilde{}} \FunctionTok{I}\NormalTok{(Temperature }\SpecialCharTok{{-}}\NormalTok{ Thermal\_Tolerance)}
\NormalTok{)}

\CommentTok{\# 2. Fit the multi{-}scale causal model}
\NormalTok{fit }\OtherTok{\textless{}{-}} \FunctionTok{because}\NormalTok{(}
    \AttributeTok{equations =}\NormalTok{ eqs,}
    \AttributeTok{data =}\NormalTok{ macroeco\_data,}
    \AttributeTok{family =} \FunctionTok{c}\NormalTok{(}\AttributeTok{Abundance =} \StringTok{"poisson"}\NormalTok{), }\CommentTok{\# Handle the non{-}Gaussian count data}
    \AttributeTok{structure =} \FunctionTok{list}\NormalTok{(}
        \AttributeTok{phylo =}\NormalTok{ anole\_tree,            }\CommentTok{\# The evolutionary scale}
        \AttributeTok{site =}\NormalTok{ my\_spatial              }\CommentTok{\# The ecological scale}
\NormalTok{    ),}
    \AttributeTok{id\_col =} \StringTok{"phylo"}                   \CommentTok{\# Map rows to the phylogenetic tree}
\NormalTok{)}
\end{Highlighting}
\end{Shaded}

By providing \texttt{structure\ =\ list(...)}, the engine automatically maps the appropriate precision matrices to the corresponding variables in the network based on the row names of the data.

In a single operation, \texttt{because} generates a full hierarchical JAGS script that simultaneously evaluates the cross-level causal paths, the phylogenetic signal (\(\lambda\)), and the spatial variance. When applied to our simulated dataset, the \texttt{because} engine successfully disentangled the dual-scale noise from the true causal signal, recovering parameter estimates that cleanly matched the true simulated values:

\begin{itemize}
\tightlist
\item
  \(\beta_{Body Mass \rightarrow Thermal Tolerance}\) (Species Level): Estimated \textbf{0.92} (\(\pm 0.11\), True value: 0.8)
\item
  \(\beta_{Elevation \rightarrow Temperature}\) (Site Level): Estimated \textbf{-1.26} (\(\pm 0.12\), True value: -1.5)
\item
  \(\beta_{Thermal Mismatch \rightarrow Abundance}\) (Observation Level): Estimated \textbf{-0.44} (\(\pm 0.09\), True value: -0.5)
\end{itemize}

The accurate recovery of the causal \(\beta\) weights across three entirely distinct levels of biological organization highlights the robustness of the fully joint hierarchical approach.

\subsubsection{5.3 The Danger of Ignoring Spatial Covariance}\label{the-danger-of-ignoring-spatial-covariance}

To empirically demonstrate the necessity of this joint architecture, we refit the exact same structural network to the simulated data, but deliberately omitted the ecological \texttt{site} matrix from the \texttt{structure} argument (simulating a standard phylogenetic path analysis that ignores spatial clustering).

While the naive model still recovered relatively accurate point estimates, the posterior standard deviations for the site and species-level effects were violently compromised. For example, the estimated standard error for the \texttt{Body\_Mass\ -\textgreater{}\ Thermal\_Tolerance} effect artificially shrank from \(\pm 0.11\) in the correct dual-scale model down to \(\pm 0.03\). By assuming that all 500 records representing observations of only 50 true species were completely independent data points, the naive model committed classic pseudoreplication. This generated an artificially narrow credibility interval and over-confident evidence for the hypothesis---a mathematical artifact that frequently leads to severe Type I error rates in empirical literature.

Because the \texttt{because} engine retains the structural topology of the user's equations, visualizing these Bayesian estimates is as simple as passing the fitted object to the native plotting function:

\begin{Shaded}
\begin{Highlighting}[]
\FunctionTok{plot\_dag}\NormalTok{(fit)}
\end{Highlighting}
\end{Shaded}

By abstracting the complex matrix algebra required for multi-scale inference, \texttt{because} empowers researchers to build and validate biological models that reflect the true complexity of the natural world without resorting to dangerous assumptions of local independence.

\section{6. Discussion: The Horizon of Joint Causal Inference}\label{discussion-the-horizon-of-joint-causal-inference}

The transition from strictly observational correlations to structured causal inference is one of the most critical methodological shifts in modern ecology and evolutionary biology (Pearl 2009; Shipley 2016). While early software packages pioneered the accessibility of these methods, the computational limitations of the era necessitated compromises, most notably the piece-wise approach to local estimation.

The \texttt{because} engine represents the natural evolution of these tools. By shifting the computational burden from matrix inversion to optimized precision matrix mapping, we can finally realize the promise of joint Bayesian estimation without sacrificing the flexibility of the local structural approach.

\subsubsection{6.1 Comparison with Existing Frameworks}\label{comparison-with-existing-frameworks}

It is essential to acknowledge the foundational role of packages like \texttt{piecewiseSEM} (Lefcheck 2016), which fundamentally democratized causal inference for ecologists. The piece-wise approach (evaluating \(X \rightarrow Y\) independently from \(Z \rightarrow X\)) is computationally lightweight and highly intuitive.

However, the assumption of local independence can fail catastrophically in the presence of hierarchical data, phylogenetic non-independence, or severe missingness. When a researcher uses \texttt{piecewiseSEM} on phylogenetic data, they must either assume the phylogenetic signal is identical across all nodes (an evolutionary impossibility) or drop the covariance structure entirely to fit the model.

In contrast, the joint Bayesian approach native to \texttt{because} evaluates the entire causal graph simultaneously. This means that uncertainty percolates mathematically through the network. If the path \(X \rightarrow M\) is highly uncertain due to missing data or phylogenetic noise, the model correctly penalizes the confidence of the downstream \(M \rightarrow Y\) path. This results in wider, more conservative credibility intervals that accurately reflect our true state of biological ignorance---a vital safeguard against Type I errors in complex systems.

\subsubsection{6.2 The Future of Extensible Inference}\label{the-future-of-extensible-inference}

The ``Mechanism-First'' design philosophy of \texttt{because} ensures that the user is focused on the biology, not the Bayesian syntax. As demonstrated via the \texttt{because.phybase} extension, the core engine was built specifically to accommodate the domain-specific nuances of evolutionary biology.

The immediate horizon for the package involves expanding this ecosystem of extensions. Development is already underway for modules like \texttt{because.occupancy}, which will seamlessly integrate detection probability and zero-inflation directly into causal networks without requiring the user to master the complexities of N-mixture models (Royle 2004). Furthermore, the abstract structure of the JAGS templates allows for rapid integration of Bayesian Model Averaging (BMA; Hoeting et al. (1999)) across topological uncertainty and adaptive priors based on formalized expert elicitation.

Ultimately, the goal of \texttt{because} is to remove the methodological friction between a biologist's hypothesis and the mathematical reality of the natural world. By automating the generation of complex, multi-scale hierarchical Bayesian models, we empower researchers to ask ``Why?'' with unprecedented rigor.

\begin{center}\rule{0.5\linewidth}{0.5pt}\end{center}

\section{Executive Summary}\label{executive-summary}

\textbf{``Because we can: Joint Causal Inference across Evolutionary and Ecological Scales''} introduces the \texttt{because} R package, a next-generation analytical engine designed to unify modern causal inference with the complex realities of biological data.

While existing software popularized Structural Equation Modeling (SEM) in ecology, they typically rely on ``piece-wise'' local estimation to skirt computational bottlenecks, leading to severe limitations when handling hierarchical cross-classification, deep-time phylogenetic covariance, and non-Gaussian data.

\texttt{because} solves this by automating the translation of intuitive \texttt{R} formulas into optimized, joint Bayesian hierarchical networks executed via JAGS. By reframing dense covariance matrices into sparse precision matrices, the engine achieves a computational paradigm shift from \(O(N^3)\) to \(O(N^2)\), enabling the simultaneous assessment of entire causal networks even on large datasets.

Key innovations include:
1. \textbf{Native Dual-Scale Integration:} The ability to simultaneously incorporate multiple, additive precision matrices (e.g., phylogenetic trees and spatial distance grids) to explicitly disentangle evolutionary inertia from concurrent ecological mechanisms.
2. \textbf{Bayesian Imputation:} Treating missing data not as an exclusionary error, but as a latent variable to be probabilistically estimated within the joint posterior.
3. \textbf{Modular Extensibility:} A design philosophy that allows domain specialists to inject custom likelihoods and covariance structures without writing a single line of MCMC code.

Through a simulated dual-scale case study, we demonstrate that \texttt{because} successfully recovers true causal effects while precisely accounting for both phylogenetic signal and spatial autocorrelation. \texttt{because} stands as a bridge between the elegance of directed acyclic graphs and the messy reality of the natural world, allowing researchers to formally test mechanistic hypotheses with robust, transparent Bayesian uncertainty.

\protect\phantomsection\label{refs}
\begin{CSLReferences}{1}{1}
\bibitem[\citeproctext]{ref-dudney2025}
Dudney, Joan, Laura E Dee, Robert Heilmayr, Jarrett Byrnes, and Katherine Siegel. 2025. {``A Causal Inference Framework for Climate Change Attribution in Ecology.''} \emph{EcoEvoRxiv}.

\bibitem[\citeproctext]{ref-felsenstein1985}
Felsenstein, Joseph. 1985. {``Phylogenies and the Comparative Method.''} \emph{The American Naturalist} 125 (1): 1--15.

\bibitem[\citeproctext]{ref-hadfield2010}
Hadfield, Jarrod D. 2010. {``MCMC Methods for Multi-Response Generalized Linear Mixed Models: The MCMCglmm r Package.''} \emph{Journal of Statistical Software} 33 (2): 1--22.

\bibitem[\citeproctext]{ref-vonhardenberg2025}
Hardenberg, Achaz von, and Alejandro Gonzalez-Voyer. 2025. {``Phylogenetic Bayesian Structural Equation Models (PhyBaSE).''} \emph{In Press/Submitted}.

\bibitem[\citeproctext]{ref-henderson1950}
Henderson, Charles R. 1950. {``Estimation of Genetic Parameters.''} \emph{Annals of Mathematical Statistics}, 309--10.

\bibitem[\citeproctext]{ref-hoeting1999}
Hoeting, Jennifer A, David Madigan, Adrian E Raftery, and Chris T Volinsky. 1999. {``Bayesian Model Averaging: A Tutorial.''} \emph{Statistical Science}, 382--401.

\bibitem[\citeproctext]{ref-lambert1992}
Lambert, Diane. 1992. {``Zero-Inflated Poisson Regression, with an Application to Defects in Manufacturing.''} \emph{Technometrics} 34 (1): 1--14.

\bibitem[\citeproctext]{ref-lefcheck2016}
Lefcheck, Jonathan S. 2016. {``piecewiseSEM: Piecewise Structural Equation Modeling in r for Ecology, Evolution, and Systematics.''} \emph{Methods in Ecology and Evolution} 7 (5): 573--79.

\bibitem[\citeproctext]{ref-legendre1993}
Legendre, Pierre. 1993. {``Spatial Autocorrelation: Trouble or New Paradigm?''} \emph{Ecology} 74 (6): 1659--73.

\bibitem[\citeproctext]{ref-little2019}
Little, Roderick JA, and Donald B Rubin. 2019. \emph{Statistical Analysis with Missing Data}. 3rd ed. John Wiley \& Sons.

\bibitem[\citeproctext]{ref-pearl2009}
Pearl, Judea. 2009. \emph{Causality: Models, Reasoning, and Inference}. 2nd ed. Cambridge University Press.

\bibitem[\citeproctext]{ref-plummer2003}
Plummer, Martyn. 2003. {``JAGS: A Program for Analysis of Bayesian Graphical Models Using Gibbs Sampling.''} \emph{Proceedings of the 3rd International Workshop on Distributed Statistical Computing} (Vienna, Austria) 124 (125.10): 1--10.

\bibitem[\citeproctext]{ref-royle2004}
Royle, J Andrew. 2004. {``N-Mixture Models for Estimating Population Size from Spatially Replicated Counts.''} \emph{Biometrics} 60 (1): 108--15.

\bibitem[\citeproctext]{ref-rue2009}
Rue, Håvard, Sara Martino, and Nicolas Chopin. 2009. {``Approximate Bayesian Inference for Latent Gaussian Models by Using Integrated Nested Laplace Approximations.''} \emph{Journal of the Royal Statistical Society: Series B (Statistical Methodology)} 71 (2): 319--92.

\bibitem[\citeproctext]{ref-shipley2000_dsep}
Shipley, Bill. 2000. {``A New Inferential Test for Path Models Based on Directed Acyclic Graphs.''} \emph{Structural Equation Modeling} 7 (2): 206--18.

\bibitem[\citeproctext]{ref-shipley2013}
Shipley, Bill. 2013. {``The AIC Model Selection Method Applied to Path Analytic Models Compared Using a d-Separation Test.''} \emph{Ecology} 94 (3): 560--64.

\bibitem[\citeproctext]{ref-shipley2016}
Shipley, Bill. 2016. \emph{Cause and Correlation in Biology: A User's Guide to Path Analysis, Structural Equations and Causal Inference with r}. 2nd ed. Cambridge University Press.

\bibitem[\citeproctext]{ref-siegel2025}
Siegel, Katherine J, and Laura E Dee. 2025. {``Foundations and Future Directions for Causal Inference in Ecological Research.''} \emph{EcoEvoRxiv}.

\end{CSLReferences}

\end{document}
